% Generated draft: requires mathematical and source review.
\documentclass[11pt]{article}
\usepackage[margin=25mm]{geometry}
\usepackage{fontspec}
\setmainfont{Latin Modern Roman}
\usepackage{amsmath,amssymb,mathtools,unicode-math}
\setmathfont{Latin Modern Math}
\usepackage{enumitem,needspace,xurl,hyperref}
\setlength{\parindent}{0pt}
\setlength{\parskip}{5pt}
\setlength{\emergencystretch}{3em}
\newcommand{\N}{\mathbb{N}}
\newcommand{\Z}{\mathbb{Z}}
\newcommand{\Q}{\mathbb{Q}}
\newcommand{\R}{\mathbb{R}}
\newcommand{\C}{\mathbb{C}}
\newcommand{\F}{\mathbb{F}}
\newcommand{\A}{\mathbb{A}}
\newcommand{\PP}{\mathbb{P}}
\newcommand{\E}{\mathbb{E}}
\newcommand{\eps}{\varepsilon}
\newcommand{\dd}{\,\mathrm{d}}
\begin{document}
\section*{UnsolvedMath ID: 30002323}
\textbf{Source problem code:} catalogue entry
\subsection{Definitions and mathematical statement}
Write \([n]=\{1,\ldots,n\}\), and let \(S_n\) be the group of all permutations of \([n]\). Under the pointwise-stabilizer interpretation, for \(A\subseteq[n]\) define
\[
H_A=\{\sigma\in S_n:\sigma(a)=a\text{ for every }a\in A\}.
\]
For \(0\leq k\leq n\), a \(k\)-coset is a left coset \(gH_A\), where \(g\in S_n\) and \(|A|=k\); the set \(A\) may vary between parts. Equivalently, it is a set
\[
C(A,f)=\{\sigma\in S_n:\sigma|_A=f\},
\]
where \(|A|=k\) and \(f:A\to[n]\) is injective. In particular, the unique \(0\)-coset is \(S_n\).

A partition of \(S_n\) is a family of pairwise disjoint nonempty subsets whose union is \(S_n\). A partition \(\mathcal P\) refines a partition \(\mathcal Q\) if every member of \(\mathcal P\) is contained in some member of \(\mathcal Q\).

\textbf{Question under this interpretation.} For every integer \(n\geq1\), every integer \(t\) with \(1\leq t\leq n\), and every partition \(\mathcal P\) of \(S_n\) into \(t\)-cosets, does there exist a partition \(\mathcal Q\) of \(S_n\) into \((t-1)\)-cosets such that
\[
\forall P\in\mathcal P\quad\exists Q\in\mathcal Q\quad P\subseteq Q?
\]
\textbf{Unresolved source qualifications.} The supplied statement does not explicitly specify the range of \(n,t\), whether the stabilizer is pointwise or setwise, or whether the stabilized set may vary between parts. The formulation above makes the usual pointwise interpretation explicit, with varying stabilized sets and the natural range \(1\leq t\leq n\); these choices require confirmation against the source. The case \(t=1\) uses the convention for \(0\)-cosets stated above.
\subsection{Short English statement}
Under the pointwise-stabilizer interpretation, each part consists of all permutations with prescribed images on \(t\) distinct points. The question asks whether every such partition can be grouped into parts that each prescribe images on \(t-1\) points. The supplied record labels the problem open; its intended parameter range and stabilizer convention require review.
\subsection{Sources}
\begin{itemize}
\item Local UnsolvedMath registry: supplied complete record, permanent ID \(30002323\), exact source problem code \texttt{catalogue entry}, title \emph{Hierarchical Refinement of Symmetric-Group Coset Partitions}. The record assigns proposed year \(2013\) and supplies no named proposer.
\item Supplied report citation: \emph{Combinatorics and Probability (2014)}, Oberwolfach Reports, DOI \texttt{10.4171/owr/2013/18}, pp.\ 1087--1152. Supplied source and literature-source link, labelled ``Oberwolfach combinatorics report'': \url{https://doi.org/10.4171/owr/2013/18}. The collection describes Oberwolfach Reports as records of mathematical workshops and their open-problem discussions. The citation's year \(2014\) differs from the proposed year \(2013\) and the year component of the DOI; this discrepancy is unresolved here.
\item Supplied literature assessment, checked on \texttt{2026-08-21}: the refinement question is described as explicitly left open, with no later resolution found in the reviewed literature. This is supplied dated curation derived from the OWR collection, not an independently verified assessment; time-sensitive status and cited claims require independent verification.
\item The registry labels the statement \texttt{corrected\_verified} and reports that its mathematical content was checked against the official Oberwolfach report, with notation, context, and question form normalized without changing the source-level claim. This verification claim and the supplied \texttt{open} status are registry metadata, not independent verification by this draft.
\end{itemize}
\end{document}
